\chapter*{RESUMEN}
\addcontentsline{toc}{chapter}{RESUMEN}

Los métodos normativos de ampacidad representan el suelo alrededor de un cable enterrado mediante propiedades homogéneas equivalentes, lo que puede ocultar zonas de baja disipación cuando la conductividad térmica varía en el entorno. Esta investigación diseñó y evaluó un artefacto basado en una red neuronal informada por física (PINN). El solucionador resuelve la ecuación de calor estacionaria 2D en el conductor, todas las capas del cable y el suelo, imponiendo continuidad de temperatura y de flujo normal en las interfaces.

Se aplicó la ciencia del diseño con una referencia de elementos finitos (FEM) refinada en tres mallas para cada caso y umbrales fijados antes de la confirmación. La búsqueda ejecutó una campaña de 101 entrenamientos dividida en cuatro etapas: arquitectura y tamaño, colocaciones, tasa de aprendizaje y peso de continuidad térmica. Luego se congeló la receta seleccionada y se evaluó su confirmación con semillas nuevas en siete casos de estudio; las temperaturas FEM no se utilizaron como etiquetas de entrenamiento.

Ningún candidato de las etapas de desarrollo cumplió los criterios en todos los problemas y semillas. La receta congelada superó la puerta de la campaña explícita en 7 de 21 ejecuciones de confirmación. Dichos criterios exigen un error térmico y NRMSE de hasta 5\,\pct{}, desequilibrio energético de hasta 2\,\pct{}, salto térmico de hasta \SI{0,1}{\kelvin} y salto de flujo de hasta 1\,\pct{}. Los casos aceptados fueron el polietileno reticulado (XLPE) homogéneo y la zona seca lejana en sus tres semillas, junto con una de tres semillas en la zona seca cercana. La configuración no fue aceptada en el arreglo plano de Aras, el caso Kim en arena, el relleno ni los estratos. La referencia FEM cuantificó el efecto físico de la heterogeneidad sobre la temperatura máxima entre \SI{-3,3}{\kelvin} (relleno) y \SI{+6,0}{\kelvin} (zona seca cercana), pero la PINN no lo resolvió bajo la regla de incertidumbre preinscrita. La búsqueda de corriente límite con la formulación explícita no se ejecutó.

Se concluyó que el artefacto es aceptable únicamente dentro del dominio ensayado y que el procedimiento trazable, al conservar los rechazos, delimita con precisión su aplicabilidad. No se demostró optimalidad global, generalización sin reentrenamiento, validación experimental de campo, ampacidad en corriente alterna (AC) normativa ni evolución temporal verificada.

\noindent\textbf{Palabras clave:} cable enterrado; ecuación de calor; PINN; interfaces; colocaciones; elementos finitos; ampacidad.

\chapter*{ABSTRACT}
\addcontentsline{toc}{chapter}{ABSTRACT}

Normative ampacity methods represent the soil around a buried cable through equivalent homogeneous properties, which can hide poorly dissipating zones when thermal conductivity varies across the surroundings. This research aimed to design and evaluate an artifact based on a physics-informed neural network (PINN) that solves the steady two-dimensional heat equation in the conductor, all cable layers, and the soil, enforcing temperature and normal heat-flux continuity across interfaces.

Design science was applied using a finite element (FEM) reference refined across three meshes per case, acceptance thresholds fixed prior to confirmation, and a campaign of 101 training runs organized into four stages: architecture and size, collocation points, learning rate, and thermal continuity weight. The selected recipe was then frozen and confirmed with new random seeds across seven cases; FEM temperatures were not used as training labels.

No development candidate met the criteria across all problems and seeds. The frozen recipe satisfied them in 7 of 21 confirmation runs under the explicit campaign tolerances (thermal error and NRMSE up to 5\,\pct{}, energy imbalance up to 2\,\pct{}, interface temperature jump up to \SI{0.1}{\kelvin}, and flux jump up to 1\,\pct{}): in all three seeds for homogeneous cross-linked polyethylene (XLPE) and the distant dry zone, and in one of three seeds for the nearby dry zone. It was not accepted for Aras flat, Kim in sand, the backfill, or the discrete soil strata. The FEM reference quantified the heterogeneity effect on maximum conductor temperature between \SI{-3.3}{\kelvin} (backfill) and \SI{+6.0}{\kelvin} (nearby dry zone), yet the PINN did not resolve this effect under the pre-registered uncertainty rule. The current-limit search with the explicit formulation was not executed.

It was concluded that the artifact is acceptable only within the tested domain and that the traceable procedure, by retaining rejected runs, rigorously bounds its applicability. Global optimality, generalization without retraining, field validation, standard-compliant AC ampacity, and verified transient evolution were not demonstrated.

\noindent\textbf{Keywords:} buried cable; heat equation; PINN; interfaces; collocation; finite elements; ampacity.
